


%*****************************************
\chapter{Introduction}
\pagenumbering{arabic} %makes sure the introduction starts at page 1
%*****************************************
\hint{This chapter should motivate the thesis, provide a clear description of the problem to be solved, and describe the major contributions of this thesis.}


% \section{Overview: Writing Your Thesis}

% \hint{
% This document provides only a general overview. 
% Although you should stick to the common structure (Intro, Background, Related Work, Design, Implementation, Evaluation, Conclusion),
% you have to adapt the sections and chapters to your individual thesis! 
% The sections in this template are not fixed and they can be (re-)moved, renamed, others can be added, etc.
% }

% \subsection{General Tips}

% \hint{
% The following list should provide you with some tips for writing your thesis with LaTeX.
% }

% \begin{itemize}
% \item 
% 	We strongly recommend that you use a LaTeX Editor like TeXMaker or TeXStudio.
%     Alternatively, you can use the university's online \LaTeX\~editor, ShareLaTeX: \url{https://www.hrz.tu-darmstadt.de/services/it_services/sharelatex_hrz/index.de.jsp}.
% \item 
%     These editors usually allow to define a file as "master file", which is then compiled every time, no matter in which file you are currently in. 
% 	In TeXMaker, it's found in the menu bar under 'options'.

% \item 
% 	To keep an overview of long paragraphs, you can put every sentence in a separate line. 
% 	A linebreak is produced with leaving a blank line (do not use \textbackslash\textbackslash or commands like \textbackslash newline).

% \item
% 	Put every chapter in its own .tex file and use the input command in the main file (main.tex) to add the chapter to your thesis. 
% 	We already did this to provide you some examples.
	
% \item
% 	Furthermore, we added examples for figures, tables, and equations, look them up. 
	
% \item
% 	You can also use comments to structure your text, using a \% sign in front of a line.
% 	% this is a comment
	
% \item	
% 	Write ToDo's or other comments that are put in your document using the ToDo notes. \todo[inline]{Note!} 
% 	You can find them in the main file (main.tex) under \enquote{Notes and Annotations}.
% 	\todo{this is a simple inline note}

% \item 
% 	Did you recognize the use of the \textbackslash enquote command to make nice quotations? 
% 	Very helpful if you \enquote{use direct citations}.
	
% \item 
% 	To prevent a linebreak between your text and a quote, use a tilde ($\sim$) as an unbreakable space before your cite command. % text~\cite 
	
% \item 
% 	\emph{Google is your friend!}
% 	Most problems you have were already encountered by others and can be easily found with a solution in the internet, 
% 	e.g., on \url{https://tex.stackexchange.com/} or just search for your problem (best in English). 

% \item 
%     We recommend using a Citation Tool like Citavi for citing in LaTex as well as helping you keep track of all of your citations. There are other Citation tools available. 
 
	
% \end{itemize}

%%
%%
% \subsection{Styling, Conventions, and more}

% \begin{itemize}
% \item 
% 	All headlines for chapters, sections, subsections, etc. are capitalized using the Title Case formatting convention.
% 	As given by the IEEE:
% 	\enquote{The page titles of all IEEE digital pages should be in title case. 
% 	All words four letters or more get capitalized, as do any nouns, pronouns, verbs, or adjectives. 
% 	Any hyphenated word will require a capital letter after the hyphen in titles (for instance, Pre-University).}\footnote{\url{https://brand-experience.ieee.org/guidelines/digital/style-guide/content/}}

% \item
% 	All captions of tables and figures are always finished with a \enquote{.} (full stop).
	
% \item 
% 	Tables, figures, and equations should be tagged with a label, such as \textbackslash label\{fig:example1\}.
% 	Use \textbackslash autoref{fig:example1} to create a label, e.g. \autoref{fig:example1}, that directly links to your figure.
% 	You can also do it manually using 'Figure$\sim$\textbackslash ref\{fig:example\}'.

% \item
% 	Capitalization in your text: it is this chapter, the table, a figure, the equation --- but Table 3, Figure~5, Equation~7, Chapter 5, Section 1.1, etc.
	
% \item 
% 	To get your bibliography in the end also correct, you may need to adapt the titles of references to get the correct capitalization.
% 	The easy way is to put all your titles in double brackets, like \enquote{title = \{\{Big Omicron and Big Omega and Big Theta\}\}}~\cite{knuth:1976}.
	
	
% \end{itemize}




\section{Motivation}
%What is the motivation for doing research in this area?

Reliable physiological monitoring is crucial for health and safety in everyday environments such as vehicles and workplaces. While current non-invasive measurement techniques---such as thermographic cameras and Photoplethysmography Imaging (PPGI)---provide a comfortable alternative to invasive sensors, they are highly susceptible to environmental noise and hardware failure (e.g., sun glare on a dashboard camera). 

To extract robust parameters despite these artifacts, the overarching research project evaluates physiological reactions using Sensor Fusion. However, effectively fusing drastically different medical signals remains computationally challenging in unconstrained environments.

Existing multimodal AI models lack the robust cross-modal predictive coding necessary to maintain accurate signal estimation when a sensor stream completely fails in a real-world, dynamic environment. To ensure high reliability, an architecture is required that can dynamically reconstruct missing data streams by learning deep correlations between thermal, visual, and physiological signals. 



\section{Problem Statement and Contribution}
%What is the problem that should be solved with this thesis?

\begin{itemize}
    \item \textbf{Main RQ:} To what extent can cross-modal predictive coding in Masked Autoencoders successfully reconstruct missing multimodal sensor signals during targeted sensor failures, and how much information is lost?
    
    \item \textbf{Sub-RQ 1:} How can 1D continuous physiological time-series data (e.g., EDA and Heart Rate) and dynamic 2D videos (RGB and Thermal face) be effectively fused into a unified Vision Transformer encoder?

    \item \textbf{Sub-RQ 2:} What is the quantitative impact of applying temporal tube masking and extremely high masking ratios (90\%+) on the reconstruction of physiological indicators from noisy video data?
\end{itemize}

\section{Outline}
How is the rest of this thesis structured?

